\chapter{Interface et navigation}
\label{ch:interface}

\section{Coque applicative}

Une fois connecté, toutes les pages métier partagent la même coque~:

\begin{description}[style=nextline, leftmargin=2.2cm]
  \item[Barre latérale] Menu principal compact ou étendu. Ordre~: \menu{Agenda}, \menu{Écoles} (libellé selon le profil), \menu{Trésorerie}, puis \menu{Référentiel} (Programmes · Groupes). Le logo ramène à \texttt{/home}.
  \item[Bandeau supérieur] Titre de la page courante, fil d'Ariane, bascule \textbf{Édition}/\textbf{Consultation}, bascule thème.
  \item[Zone de contenu] Onglets de module et tableaux selon le contexte.
\end{description}

La sidebar démarre en mode \textbf{compact}~; cliquez sur le chevron pour l'élargir. L'état est mémorisé dans le navigateur (\texttt{localStorage}).

\section{Onglets par module}

Chaque grand module propose des sous-onglets~:

\begin{table}[htbp]
  \centering
  \caption{Onglets des modules principaux}
  \begin{tabular}{@{}lp{9cm}@{}}
    \toprule
    \textbf{Module} & \textbf{Onglets} \\
    \midrule
    Agenda          & Planning · Calendrier (import ICS, admin) --- vues \textbf{mois} / \textbf{semaine} \\
    Trésorerie      & Synthèse · Factures · Banque · Dépenses \\
    Référentiel     & Programmes · Groupes \\
    Paramètres      & Entreprise · Profil utilisateur \\
    \bottomrule
  \end{tabular}
\end{table}

Les actions \menu{Créer facture} et \menu{Créer dépense} sont des boutons dans la page, pas des onglets.
La préparation de facturation reste sur la fiche école (\texttt{\#billing}), pas dans la sidebar.

\section{Modes Édition et Consultation}

Le bandeau permet de basculer entre deux modes globaux~:

\begin{itemize}[leftmargin=*]
  \item \textbf{Consultation}~: lecture seule, idéal pour les reportings.
  \item \textbf{Édition}~: formulaires, boutons d'action, création et modification des enregistrements.
\end{itemize}

Vérifiez le mode actif avant de planifier une session ou de préparer une facture.

\section{Plan de charge vs liste des écoles}

Deux écrans proches mais complémentaires~:

\begin{itemize}[leftmargin=*]
  \item \texttt{/home} — \textbf{Liste des écoles}~: porte d'entrée avec \emph{facturé TTC} et \emph{non facturé TTC} par établissement, graphique de répartition (voir figure~\ref{fig:04-home.png}).
  \item \texttt{/school/dashboard} — \textbf{Plan de charge}~: KPIs annuels et tableaux de cours détaillés.
\end{itemize}

\section{Paramètres entreprise}

Accessible via \menu{Paramètres}, l'onglet \menu{Entreprise} regroupe~:

\begin{itemize}[leftmargin=*]
  \item Identité de la société (raison sociale, adresse, SIRET, TVA).
  \item Contexte métier et libellés associés.
  \item Compte bancaire de facturation par défaut.
  \item Gestion des utilisateurs (réservé à l'administrateur).
\end{itemize}

\section{Raccourcis de navigation}

\begin{workflowstep}[Retour accueil]
  Cliquez sur le logo ou sur \menu{Écoles} dans la barre latérale → \texttt{/home}.
\end{workflowstep}

\begin{workflowstep}[Ouvrir une école]
  Depuis la liste, cliquez sur le nom de l'établissement → fiche détaillée à onglets (\texttt{?panel=})~: cours + facturation, groupes, détails, documents.
\end{workflowstep}

\begin{workflowstep}[Accéder à l'agenda]
  \menu{Agenda} dans la sidebar → \texttt{/planning}.
\end{workflowstep}

\begin{workflowstep}[Accéder à la trésorerie]
  \menu{Trésorerie} → \texttt{/treasury} (synthèse mensuelle).
\end{workflowstep}

\screenshot[keepaspectratio]{07-programs.png}{Module Programmes — création des offres globales à l'entreprise}

\screenshot[keepaspectratio]{08-groups.png}{Module Groupes — liaison cours / participants}
